\SectionStart
\section{数据处理与评价指标}
\subsection{数据接口与单位}
\WritingContract
  {建立可复现的数据入口和统一计算尺度。}
  {说明实际使用的字段、样本范围、时间与空间粒度、单位换算、表间关联键及各问共享接口。}
  {字段和范围来自正式附件及运行清单，换算关系能够在数据准备代码中定位。}
  {不得复制整张原始数据表，不得使用未声明的外部数据，不得隐去会改变结果的筛选。}
\subsection{缺失值、异常值与特征构造}
\WritingContract
  {说明哪些处理改变了模型输入，以及这些处理为何合理。}
  {报告缺失和异常的识别口径、处理比例与理由；说明标准化、聚合或特征构造的公式和单位影响。}
  {所有处理由数据剖析统计和可复现脚本支撑，处理前后样本量或守恒关系可核对。}
  {不得为了篇幅罗列未实际执行的方法，不得把绘图美化当作数据处理，不得无依据删除异常值。}
\subsection{评价指标与共同口径}
\WritingContract
  {统一各问结果比较、误差评价和可行性判断的尺度。}
  {定义指标公式、方向、单位、聚合方式、基线和必要的守恒关系；题型专属指标可在对应问补充。}
  {指标定义与实验输出字段、表格列和冻结主张一致，参数来源可追溯。}
  {不得只写指标名称，不得混用平均值与总量、比例与百分点，也不得事后选择有利指标。}
